\section{Discussion}

\paragraph{Diagnostic signals are not control knobs.}
A reproducible subspace signature can help localize and compare adaptation, yet the tested interventions show that manipulating the most obvious geometric quantities is not enough to control generated behavior. This discourages a common inference from correlation to mechanism and motivates evaluating safety interventions at the sequence-behavior level.

\paragraph{Behavioral supervision has both amount and timing.}
The historical \constanttwenty{} contrast shows that sustained supervision can reduce rebound, but it mixes timing with cumulative dose. The equal-dose experiment is designed to identify the temporal component. If the temporal ordering reproduces, late effective contribution becomes a stronger design constraint. If it disappears, the scientifically correct conclusion is that cumulative dose explains the historical advantage.

\paragraph{Safety persistence is not response quality.}
Persistent supervision may hold the model on a narrow safe response policy. This can be preferable to harmful compliance, but it is not a complete outcome. Prompt binding, safe actionability, detail, naturalness, and benign utility must be measured alongside unsafe rate.

\paragraph{Implications for future methods.}
A future method should preserve the behaviorally active sequence-level signal, maintain sufficient late contribution, use reliable prompt-aware targets, and verify finite updates rather than relying only on local geometry or proxy movement. These are design requirements derived from Paper A; they are not method-validation results.

\paragraph{Value of negative results.}
The spectral, margin, token-proxy, target-generation, and measurement failures narrow the viable hypothesis space. Their value depends on bounded wording and faithful implementation checks, not on claiming universal impossibility.
